\documentclass[11pt,letterpaper]{article}
\usepackage[margin=1in]{geometry}
\usepackage{amsmath,amssymb}
\usepackage{graphicx}
\usepackage{booktabs}
\usepackage{hyperref}
\usepackage{fancyhdr}
\usepackage{titlesec}

% Header and footer setup
\pagestyle{fancy}
\fancyhf{}
\lhead{\textbf{Westchester County Data Platform} - Methodology Report}
\rhead{Page \thepage}
\renewcommand{\headrulewidth}{0.4pt}

% Title formatting
\titleformat{\section}{\Large\bfseries}{\thesection}{1em}{}
\titleformat{\subsection}{\large\bfseries}{\thesubsection}{1em}{}

% Document metadata
\title{\textbf{Westchester County Data Platform}\\Methodology Report}
\author{Arcanum Research Initiative}
\date{October 2025}

\begin{document}

\maketitle

\tableofcontents
\newpage

\section{Executive Summary}

\textbf{Project Purpose:} The Westchester County Data Platform provides comprehensive analysis of demographic, economic, transit, and municipal data for Westchester County, New York through integrated open data sources, interactive mapping, and professional-grade dashboards.

\textbf{Methodology Overview:} Multi-source data collection from U.S. Census Bureau, Metro-North Railroad GTFS feeds, and NY State Open Data portal, processed into Druck-compliant Excel files and interactive web visualizations.

\textbf{Key Findings:} Successfully integrated 56 Metro-North stations, 241 census tracts, and county-level demographics for comprehensive geographic and statistical analysis.

\textbf{Deliverables:} 4 Excel datasets, interactive React/TypeScript web platform, FastAPI backend, and comprehensive documentation.

\section{Research Objectives}

\subsection{Primary Research Questions}
\begin{itemize}
    \item What are the current demographic characteristics of Westchester County?
    \item How accessible is public transit (Metro-North) across the county?
    \item What geographic and economic patterns exist across municipalities?
    \item How can open data sources be integrated for evidence-based policy analysis?
\end{itemize}

\subsection{Success Criteria}
\begin{itemize}
    \item Complete data collection from all primary sources (Census, GTFS, NY State)
    \item Generation of Druck-compliant Excel files (ONE sheet per file)
    \item Operational web platform with interactive dashboards
    \item Professional documentation suitable for government stakeholders
\end{itemize}

\section{Methodology}

\subsection{Data Sources}

\begin{itemize}
    \item \textbf{U.S. Census Bureau API:} American Community Survey (ACS) 5-Year Estimates (2022)
    \begin{itemize}
        \item County-level demographics (population, income, housing, employment)
        \item Census tract-level data (241 tracts)
        \item Municipality-level aggregations (1,285 places in NY State, filtered to Westchester)
        \item API endpoint: \texttt{https://api.census.gov/data/2022/acs/acs5}
        \item No API key required (recommended for higher rate limits)
    \end{itemize}
    
    \item \textbf{Metro-North Railroad GTFS:} General Transit Feed Specification
    \begin{itemize}
        \item Station locations, routes, and accessibility data
        \item Downloaded from: \texttt{http://web.mta.info/developers/data/mnr/google\_transit.zip}
        \item Geographic filtering: Lat 40.9-41.4, Lon -74.0 to -73.5
        \item Result: 56 stations in Westchester County
    \end{itemize}
    
    \item \textbf{NY State Open Data (Socrata):} data.ny.gov platform
    \begin{itemize}
        \item Crime statistics (1,611 records for Westchester)
        \item Health facilities (330 facilities)
        \item Property assessment data (pending integration)
        \item App token recommended for API access
    \end{itemize}
\end{itemize}

\subsection{Analytical Framework}

\textbf{Core Approach:} Geographic Information Systems (GIS) analysis combined with statistical demographic analysis.

\textbf{Transit Accessibility Model:}
\begin{equation}
    \text{Coverage Area} = \pi r^2 \times N_{stations}
\end{equation}
where $r = 1$ mile (walkable radius) and $N_{stations} = 56$

\textbf{Implementation Steps:}
\begin{enumerate}
    \item Data Collection: Automated API requests to Census, GTFS download, Socrata queries
    \item Geographic Filtering: Bounding box filter for Westchester County coordinates
    \item Data Transformation: JSON/CSV conversion, column standardization
    \item Excel Generation: Druck-compliant formatting (one sheet per file)
    \item Visualization: React/TypeScript dashboards with Recharts and Leaflet mapping
\end{enumerate}

\subsection{Quality Assurance}

\textbf{Validation Procedures:}
\begin{itemize}
    \item Automated Excel validation (ONE sheet per file requirement)
    \item Geographic coordinate verification (within Westchester bounds)
    \item Data completeness checks (non-null critical fields)
    \item Cross-reference with official county statistics
\end{itemize}

\textbf{Accuracy Metrics:} 100\% compliance with Druck standards, 0 critical violations in Excel validation.

\section{Implementation Details}

\subsection{Software and Tools}
\begin{itemize}
    \item \textbf{Primary Analysis:} Python 3.11+ with pandas, geopandas, requests
    \item \textbf{Data Processing:} Custom processors for demographics and transit
    \item \textbf{Excel Generation:} openpyxl with Druck formatting standards
    \item \textbf{Backend API:} FastAPI with 8 operational endpoints
    \item \textbf{Frontend:} React 18, TypeScript, Vite, Tailwind CSS
    \item \textbf{Mapping:} Leaflet for interactive geographic visualization
    \item \textbf{Charts:} Recharts for statistical visualizations
\end{itemize}

\subsection{Technical Architecture}
\begin{itemize}
    \item \textbf{Data Layer:} Raw data in \texttt{Technical/data/raw/}, processed in \texttt{Technical/data/processed/}
    \item \textbf{API Layer:} FastAPI REST endpoints serving JSON data
    \item \textbf{Frontend Layer:} React SPA with client-side routing
    \item \textbf{Output Layer:} Excel files in \texttt{Output/Data/Results/}, PDFs in \texttt{Output/PDFs/}
\end{itemize}

\section{Data Processing Pipeline}

\subsection{Step 1: Data Acquisition}
\begin{verbatim}
python Technical/scripts/download_all_data.py
\end{verbatim}
Downloads all data sources, handles API authentication, geographic filtering

\subsection{Step 2: Excel Generation}
\begin{verbatim}
python Technical/src/processors/generate_all_excel.py
\end{verbatim}
Processes raw data into Druck-compliant Excel files with proper formatting

\subsection{Step 3: Validation}
\begin{verbatim}
python Technical/scripts/validate_excel.py
\end{verbatim}
Verifies ONE sheet per file, validates column names, checks file integrity

\section{Limitations and Future Work}

\subsection{Current Limitations}
\begin{itemize}
    \item Municipality filtering requires refinement (captured 1,285 NY places, needs Westchester-specific list)
    \item Property tax data partially integrated (parcel-level analysis pending)
    \item Historical trend analysis limited to single-year snapshot (2022)
    \item Transit ridership data not included in current GTFS import
\end{itemize}

\subsection{Future Enhancements}
\begin{itemize}
    \item Multi-year time series analysis
    \item Advanced transit coverage heat mapping
    \item Integration with county budget systems
    \item Real-time data updates via scheduled jobs
    \item PostgreSQL + PostGIS for spatial queries
\end{itemize}

\section{References}

\begin{itemize}
    \item U.S. Census Bureau. (2022). \textit{American Community Survey 5-Year Estimates}. Retrieved from https://www.census.gov/programs-surveys/acs
    \item Metropolitan Transportation Authority. (2025). \textit{Metro-North Railroad GTFS Data}. Retrieved from http://web.mta.info/developers/
    \item New York State. (2025). \textit{Open Data Portal}. Retrieved from https://data.ny.gov
    \item Westchester County GIS. (2025). \textit{Tax Parcel Data}. Retrieved from Westchester County Open Data
\end{itemize}

\section{Appendices}

\subsection{Appendix A: API Endpoints}
\begin{itemize}
    \item \texttt{GET /api/municipalities} - List of county municipalities
    \item \texttt{GET /api/transit/stations} - Metro-North station data
    \item \texttt{GET /api/demographics/county} - County-level demographics
    \item \texttt{GET /api/demographics/tracts} - Census tract data
    \item \texttt{GET /api/demographics/municipalities} - Municipality demographics
    \item \texttt{GET /api/stats} - Platform statistics and data availability
\end{itemize}

\subsection{Appendix B: Excel File Inventory}
\begin{enumerate}
    \item \texttt{[2025.10.13] westchester\_county\_demographics.xlsx}
    \item \texttt{[2025.10.13] westchester\_census\_tracts.xlsx}
    \item \texttt{[2025.10.13] westchester\_metro\_north\_stations.xlsx}
    \item \texttt{[2025.10.13] westchester\_transit\_accessibility.xlsx}
\end{enumerate}

\subsection{Appendix C: Census Variables}
\begin{table}[h]
\centering
\begin{tabular}{ll}
\toprule
\textbf{Variable Code} & \textbf{Description} \\
\midrule
B01003\_001E & Total Population \\
B01002\_001E & Median Age \\
B19013\_001E & Median Household Income \\
B25077\_001E & Median Home Value \\
B23025\_005E & Unemployed Population \\
\bottomrule
\end{tabular}
\caption{Key Census ACS Variables Used}
\end{table}

\end{document}

